% !TEX root = ../tfg_etsiinf_NombreApellidos.tex
\chapter{Resultados} \label{chp:resultados}

Este capítulo presenta los resultados obtenidos mediante la estrategia de pruebas
definida previamente. La exposición distingue entre resultados de infraestructura
y resultados de aprendizaje. Esta distinción es esencial en el presente trabajo:
la capacidad de lanzar PPO y generar métricas no implica por sí misma que se haya
obtenido una política robusta, especialmente cuando los experimentos revelan
problemas de reinicio, accionabilidad o coherencia física del simulador.

\section{Validación funcional de la infraestructura} \label{sct:resultados:validacion}

Las primeras pruebas confirman que la infraestructura básica es operativa. El
entorno puede instanciarse como entorno Gymnasium, conectarse con Aerostack2,
armar el UAV, activar el modo \emph{offboard}, ejecutar despegue, enviar comandos
de velocidad y registrar métricas mediante los mecanismos de monitorización. Esta
validación funcional demuestra que la integración entre Gymnasium,
Stable-Baselines3, ROS 2 y Aerostack2 es viable.

No obstante, los resultados también muestran que la validez de un entrenamiento
depende de condiciones adicionales. En las primeras ejecuciones se observaron
éxitos triviales cuando el objetivo estaba demasiado cerca de la pose de
despegue, ausencia de escalares de TensorBoard cuando PPO no alcanzaba suficientes
pasos de \emph{rollout}, y bucles de episodios fuera de límites cuando el reset no
garantizaba un estado inicial seguro. Las iteraciones posteriores permitieron
certificar el reinicio mediante el servicio del simulador y ejecutar
entrenamientos más largos, pero también revelaron que las salidas de límites y la
degradación acumulada por reinicios de servicio podían dominar el comportamiento
observado. Estos resultados condujeron a separar el problema en dos niveles:
primero validar el ciclo físico de episodio y después analizar el aprendizaje.

\section{Matriz de experimentos PPO} \label{sct:resultados:matriz_experimentos}

El Cuadro~\ref{tab:resultados:experimentos} resume la evolución de los
experimentos principales. La tabla no debe interpretarse como una búsqueda lineal
de hiperparámetros, sino como una secuencia de validación en la que cada ejecución
revela una propiedad de la infraestructura.

\begin{table}[H]
\centering
\scriptsize
\setlength{\tabcolsep}{2.5pt}
\renewcommand{\arraystretch}{1.12}
\begin{tabular}{p{0.10\textwidth} p{0.21\textwidth} p{0.25\textwidth} p{0.34\textwidth}}
\hline
\textbf{Experimento} & \textbf{Objetivo} & \textbf{Resultado observado} & \textbf{Interpretación técnica} \\
\hline
Exp001 & Smoke PPO con configuración inicial. & Episodios de éxito en un paso y métricas insuficientes. & El objetivo estaba demasiado cerca del estado inicial; no era una tarea de navegación útil. \\
Exp002 & Objetivo desplazado y TensorBoard en prueba corta. & Primeros episodios útiles, después bucles de salida de límites. & El pipeline generaba artefactos, pero el reset no garantizaba inicio seguro tras OOB. \\
Exp003 & Introducir \texttt{fixed\_start\_pose}. & Entrenamiento corto completado con monitor, TensorBoard y \emph{checkpoints}. & El reset determinista resolvió el bucle inmediato de OOB en una prueba acotada. \\
Exp004 & Target más lejano y marcador RViz. & Se mejoró la observabilidad del objetivo y de métricas. & RViz y marcadores ayudan a validar cualitativamente la escena, no sustituyen métricas. \\
Exp005 & Target estable a 2 m y reset en vuelo. & Primeras evidencias de éxito real hacia un punto, con incidencias de reset. & El reset en vuelo evitó bloqueos de aterrizaje, pero la robustez seguía siendo limitada. \\
Exp006 & Mejorar estabilidad y evitar fallback problemático. & Se detectó bloqueo asociado a \texttt{go\_to}. & El retorno a pose inicial debía evitar comportamientos bloqueantes. \\
Exp007 & Gate largo tras eliminar \texttt{go\_to}. & Fallo acotado por timeout de reset por velocidad. & El entrenamiento ya no quedaba bloqueado, pero el reset seguía sin ser fiable para runs largos. \\
Exp008 & Reset más robusto con target de 2 m. & Persisten timeouts de reset. & La dificultad seguía siendo alta para depurar aprendizaje si el reset fallaba. \\
Exp008a & Curriculum de 0,5 m y diagnóstico ampliado. & Mejora parcial, pero aparecen bloqueos de reset y accionabilidad. & El cuello de botella pasa de PPO a la interacción reset--controlador--simulador. \\
Exp009 & Primer gate PPO sobre el stack de reset certificado. & Ejecución de 5102 pasos de monitor, 126 episodios y finalización correcta del proceso. & Se valida la infraestructura de episodio; la tasa de éxito observada, 21/126 episodios, no se interpreta todavía como política robusta. \\
Exp010 & Recompensa mínima alineada con el tutor, randomización completa e hiperparámetros revisados. & Se registran 12852 pasos y 427 episodios antes del fallo; 359 terminan fuera de límites y 65 por baja altitud. & El diseño de recompensa queda alineado, pero las fronteras no observables y el coste de teleport dominan la experiencia. \\
Exp011 & Muestreo en área segura y límite de distancia inicio--objetivo para evitar saturación de observaciones. & Se registran 21347 pasos y 637 episodios; 512 terminan fuera de límites, 121 por baja altitud y solo 3 en éxito. & La saturación inicial se corrige parcialmente, pero el problema estructural de terminaciones y reinicios repetidos persiste. \\
Exp012 & Guard de acción en límites, corredor vertical ampliado, reanudación desde \emph{checkpoint} y supervisor. & Experimento en ejecución durante la redacción de esta versión. & La hipótesis evaluada es reducir reinicios de servicio y aislar mejor el aprendizaje de PPO de los fallos de infraestructura. \\
\hline
\end{tabular}
\caption{Resumen de experimentos principales y su interpretación.}
\label{tab:resultados:experimentos}
\end{table}

\section{Resultados de entrenamiento y monitorización} \label{sct:resultados:entrenamiento}

Los experimentos muestran que la infraestructura de entrenamiento produce
artefactos analizables cuando alcanza suficiente número de pasos: ficheros de
monitor, eventos de TensorBoard, \emph{checkpoints} y modelos finales en
pruebas acotadas. En particular, las ejecuciones iniciales permiten comprobar que
TensorBoard y los monitores de episodio funcionan cuando PPO completa
\emph{rollouts} suficientes, mientras que Exp009 confirma que el stack de reset
certificado puede sostener una ejecución PPO completa de 5000 pasos. Sin embargo,
la presencia de estos artefactos no implica aprendizaje estable: debe analizarse
junto con las razones de terminación, la longitud de episodio, la distancia
final, la tasa de éxito y los diagnósticos de reset.

Las métricas agregadas de TensorBoard son especialmente útiles para interpretar
Exp010 y Exp011, porque permiten observar si la recompensa media, la longitud de
episodio o la tasa de éxito muestran una tendencia temporal. En las condiciones
registradas, la evidencia disponible no debe presentarse como una política final
aprendida, sino como una secuencia de diagnóstico: Exp010 muestra que la
formulación de recompensa e hiperparámetros ya está alineada con el criterio del
tutor, pero que la experiencia queda dominada por salidas de límites; Exp011
reduce la saturación de observaciones mediante un área de muestreo más segura,
pero no elimina el predominio de terminaciones no deseadas. Por ello, las curvas
de TensorBoard deben leerse junto con los ficheros de monitor y con la matriz de
experimentos, no de forma aislada.

La comparación cuantitativa entre Exp010 y Exp011 muestra una mejora parcial,
pero insuficiente. En Exp010 se registran 427 episodios, con 359 terminaciones
por salida de límites, 65 por baja altitud insegura y 3 éxitos. La recompensa
media por episodio es aproximadamente \(-31{,}36\), la longitud media es de
30,10 pasos y la distancia final media al objetivo es de 5,41 m. En Exp011 se
registran 637 episodios, con 512 salidas de límites, 121 terminaciones por baja
altitud, 3 éxitos y un episodio por límite de pasos. La recompensa media mejora
hasta aproximadamente \(-27{,}34\), la longitud media aumenta a 33,51 pasos y la
distancia final media baja a 3,71 m. Sin embargo, la tasa de éxito permanece por
debajo del 1\% en ambos casos y cae a cero en los tramos finales de las curvas de
TensorBoard. Por tanto, Exp011 no demuestra aprendizaje robusto, sino que reduce
algunos síntomas de Exp010 sin eliminar la causa principal: la política sigue
interactuando con límites que no forman parte explícita de la observación.

La evolución de los experimentos no permite afirmar todavía que se haya obtenido
una política PPO robusta para navegación estable. Sí permite afirmar que se ha
construido una infraestructura capaz de revelar por qué un entrenamiento no es
válido. Las ejecuciones detectaron éxitos triviales, episodios contaminados por
reinicios inválidos, bloqueos del comportamiento \texttt{go\_to}, timeouts de
reset por velocidad, comandos aceptados que no se traducían en desplazamiento
físico suficiente y degradación asociada a un número elevado de teleports del
simulador. Estos resultados desplazan la discusión desde la optimización de PPO
hacia la validez del entorno de entrenamiento.

\section{Análisis de fallos de integración} \label{sct:resultados:fallos_integracion}

El resultado más relevante desde el punto de vista de ingeniería es la
identificación del reinicio y de la accionabilidad como cuellos de botella. En
un entorno de aprendizaje por refuerzo, cada episodio debe comenzar en un estado
bien definido y cada acción debe modificar el estado de forma coherente. Cuando
el UAV no regresa a la pose inicial, cuando el controlador no mantiene la pose
tras el reset o cuando la API acepta comandos que no producen movimiento, el
entrenamiento deja de generar experiencia válida.

La solución adoptada consiste en hacer visibles esos fallos y detener la
ejecución antes de acumular datos no representativos. Esta política reduce el
número de ejecuciones largas aparentemente productivas, pero aumenta la calidad
de la evidencia experimental. En las condiciones evaluadas, el problema
principal no se encuentra en la llamada estándar a \texttt{model.learn} de
Stable-Baselines3, sino en garantizar que el entorno robótico entregue
transiciones físicamente accionables después de cada reinicio.

La incorporación del servicio de reinicio del simulador resolvió una parte
fundamental del problema, al permitir teletransportar el UAV a una pose inicial
controlada y reinicializar estado dinámico, referencias y sincronización de la
plataforma. Sin embargo, Exp010 y Exp011 muestran que un servicio de reset
correcto no basta si la tarea genera demasiadas terminaciones en fronteras que la
política no observa directamente. Esta observación justifica la introducción en
Exp012 de guards de acción en los límites de la escena: el objetivo no es ocultar
fallos de seguridad, sino impedir que un agente no entrenado convierta los
límites invisibles en una fuente dominante de reinicios costosos.

\section{Discusión de limitaciones} \label{sct:resultados:limitaciones}

Los resultados deben interpretarse con varias limitaciones. En primer lugar, las
ejecuciones disponibles combinan gates de infraestructura y entrenamientos
parciales, no una campaña final de entrenamiento cerrada. En segundo lugar, la
comparación con una configuración PID queda planteada como evaluación posterior,
ya que no sería metodológicamente correcto comparar un controlador de referencia
contra una política PPO cuyo entrenamiento estable aún depende de reducir
terminaciones espurias y reinicios costosos. En tercer lugar, los resultados
dependen del estado del simulador y de la correcta sincronización entre
plataforma, controlador y servicios de reset.

Estas limitaciones no invalidan el trabajo realizado. Al contrario, delimitan con
precisión qué parte de la infraestructura está validada y qué parte requiere
investigación adicional. La contribución experimental del trabajo reside en haber
construido una base reproducible y en haber identificado condiciones de fallo que
deben resolverse antes de extraer conclusiones sobre el rendimiento final de PPO.
